% プロコン 2026 展示ポスター A0 2 枚版の 2 枚目（AI の仕組みと評価）
% 版面は A1 で組み、poster2_a0.tex が √2 倍して A0 2 ページの poster2_a0.pdf にする。
% 評価の数値とグラフは ../../resume/poster_summary.csv・poster_curve.csv（resume/poster_data.py の出力）による
% ビルド: build_poster.bat をダブルクリック（Tectonic）
\documentclass[24pt,a1paper,portrait,margin=0mm,innermargin=12mm,
  blockverticalspace=5mm,colspace=9mm,subcolspace=8mm]{tikzposter}
% poster2_p1.tex / poster2_p2.tex（A0 2 枚版）で共有する設定。poster.tex（1 枚版）の設定を元に、文字を大きくしたもの
% 寸法・文字サイズは A1 のときの値（poster2_a0.tex で 1.414 倍して A0 にする）
\usepackage{xeCJK}
\setCJKmainfont[BoldFont=HaranoAjiGothic-Bold.otf]{HaranoAjiGothic-Regular.otf}
\setCJKsansfont[BoldFont=HaranoAjiGothic-Bold.otf]{HaranoAjiGothic-Regular.otf}
\setmainfont{texgyreheros-regular.otf}[BoldFont=texgyreheros-bold.otf]
\setsansfont{texgyreheros-regular.otf}[BoldFont=texgyreheros-bold.otf]
\xeCJKsetup{CJKspace=false}

\usepackage{amsmath,amssymb,bm}
\usepackage{enumitem}
\usepackage{graphicx}
\usepackage{array}
\usepackage{fontawesome5}
\usepackage{pgfplots}
\pgfplotsset{compat=1.18}
\usetikzlibrary{positioning,arrows.meta,calc}
\graphicspath{{../../resume/}{../../}}

\setlist{nosep,leftmargin=1.0em,itemsep=4pt}
\renewcommand{\labelitemi}{\textcolor{blue}{\textbullet}}
\setlength{\parindent}{0pt}
\setlength{\parskip}{2pt}

% ---------------------------------------------------------------- 色（DESIGN.md とロゴに合わせる）
\definecolor{navy}{HTML}{164A86}     % ロゴの文字色
\definecolor{blue}{HTML}{0093C4}     % blue-deep
\definecolor{skytint}{HTML}{E3F7FF}
\definecolor{yellow}{HTML}{FFDB4F}
\definecolor{yellowdeep}{HTML}{C99A00}
\definecolor{yellowtint}{HTML}{FFF6DA}
\definecolor{ink}{HTML}{232323}
\definecolor{ink2}{HTML}{5E584C}
\definecolor{paper}{HTML}{F1EFE8}    % ロゴの背景色
\definecolor{card}{HTML}{F6F3EA}
\definecolor{good}{HTML}{1C7A45}
\definecolor{goodtint}{HTML}{E3F4E9}
\definecolor{orange}{HTML}{E2612B}

\definecolorstyle{gemmba}{
  \colorlet{colorOne}{navy}\colorlet{colorTwo}{blue}\colorlet{colorThree}{yellow}
}{
  \colorlet{backgroundcolor}{paper}
  \colorlet{framecolor}{paper}
  \colorlet{titlefgcolor}{navy}
  \colorlet{titlebgcolor}{paper}
  \colorlet{blocktitlebgcolor}{white}
  \colorlet{blocktitlefgcolor}{navy}
  \colorlet{blockbodybgcolor}{white}
  \colorlet{blockbodyfgcolor}{ink}
  \colorlet{innerblocktitlebgcolor}{white}
  \colorlet{innerblocktitlefgcolor}{navy}
  \colorlet{innerblockbodybgcolor}{white}
  \colorlet{innerblockbodyfgcolor}{ink}
  \colorlet{notefgcolor}{ink}
  \colorlet{notebgcolor}{skytint}
  \colorlet{noteframecolor}{blue}
}
\usecolorstyle{gemmba}
\usetitlestyle{Filled}
\usebackgroundstyle{Default}

% ブロック：白いカード。見出しは中央、黄色の下線
\defineblockstyle{gemblock}{
  titlewidthscale=1, bodywidthscale=1, titlecenter,
  titleoffsetx=0pt, titleoffsety=0pt, bodyoffsetx=0mm, bodyoffsety=0mm,
  bodyverticalshift=0mm, roundedcorners=10, linewidth=0pt,
  titleinnersep=3mm, bodyinnersep=6mm
}{
  \ifBlockHasTitle
    \draw[draw=none, fill=white, rounded corners=\blockroundedcorners]
      (blockbody.south west) rectangle (blocktitle.north east);
  \else
    \draw[draw=none, fill=white, rounded corners=\blockroundedcorners]
      (blockbody.south west) rectangle (blockbody.north east);
  \fi
}
\useblockstyle{gemblock}


% ---------------------------------------------------------------- タイトル（2 枚とも同じ帯。右上にページ番号）
% \posttitle{ページの主題}{n/2}
\newcommand{\posttitle}[2]{%
  \settitle{%
    \begin{minipage}[c]{0.36\linewidth}
      \includegraphics[width=0.95\linewidth]{logo.png}
    \end{minipage}\hfill
    \begin{minipage}[c]{0.62\linewidth}
      \raggedright\color{navy}
      {\fontsize{66pt}{70pt}\selectfont\bfseries 中小企業向けのDX推進システム}\par\vspace{10pt}
      {\fontsize{36pt}{40pt}\selectfont\bfseries\color{blue}#1}\quad
      {\fontsize{26pt}{30pt}\selectfont\bfseries\color{white}\colorbox{navy}{\,#2\,}}\par\vspace{10pt}
      {\fontsize{24pt}{28pt}\selectfont\color{ink2}明石工業高等専門学校\quad 小笠原\,悠太\quad 増永\,惺真\qquad 指導教員：野村\,隼人}
    \end{minipage}}}

% ---------------------------------------------------------------- 部品
% ブロック見出し：文字幅に合わせた黄色の下線
\newcommand{\bt}[1]{\tikz[baseline=(t.base)]{%
  \node[inner sep=0, font=\fontsize{48pt}{54pt}\selectfont\bfseries, text=navy] (t) {#1};
  \fill[yellow] ([xshift=-6mm, yshift=-16pt]t.south west) rectangle ([xshift=6mm, yshift=-6pt]t.south east);}}
\newcommand{\sm}{\fontsize{20pt}{25pt}\selectfont}
\newcommand{\capt}[1]{{\centering\fontsize{20pt}{25pt}\selectfont\color{ink2}#1\par}}
\newcommand{\head}[1]{{\fontsize{32pt}{38pt}\selectfont\bfseries\color{navy}#1}\par\vspace{10pt}}
\newcommand{\hl}[1]{\textcolor{blue}{\textbf{#1}}}
\newcommand{\bl}{\fontsize{24pt}{31pt}\selectfont}   % 本文・箇条書き
\newcommand{\tb}{\fontsize{25pt}{30pt}\selectfont}   % 表
\renewcommand{\arraystretch}{1.45}
\setlist{nosep,leftmargin=1.0em,itemsep=7pt}

% 写真：img/<ファイル名> があれば表示、無ければ「ここに写真」の枠を描く
%   \photo{ファイル名}{幅}{高さ}{何の写真か}
\newcommand{\photo}[4]{%
  \IfFileExists{img/#1}{%
    \tikz{\node[inner sep=0] {\includegraphics[width=#2,height=#3,keepaspectratio]{img/#1}};}}{%
    \tikz{\draw[dashed, line width=2pt, draw=ink2!55, fill=card, rounded corners=8pt]
            (0,0) rectangle (#2,#3);
          \node[align=center, text=ink2, font=\fontsize{20pt}{25pt}\selectfont]
            at ($(0,0)!0.5!(#2,#3)$)
            {{\fontsize{36pt}{36pt}\selectfont\faCamera}\\[4pt]#4\\{\fontsize{15pt}{18pt}\selectfont img/\detokenize{#1}}};}}}
% 管理画面のスクリーンショット（細い枠つき）
\newcommand{\shot}[2]{\fcolorbox{ink2!40}{white}{\includegraphics[width=#1]{img/#2}}}

% メリット：短い見出し＋補足 1 行。項目内は詰め、項目の間を空ける
\newcommand{\merit}[2]{% 見出し, 補足
  {\bfseries\fontsize{30pt}{36pt}\selectfont #1\par}
  \vspace{3pt}
  {\fontsize{23pt}{28pt}\selectfont\color{ink2}#2\par}
  \vspace{14pt}}

% 囲み（淡い色の地のカード）\tint{地の色}{中身}
\newcommand{\tint}[2]{\coloredbox[bgcolor=#1, fgcolor=ink, framecolor=#1, roundedcorners=8, linewidth=0pt]{#2}}

\tikzset{
  box/.style={draw=navy, line width=3pt, rounded corners=8pt, fill=white,
    align=flush center, inner sep=9pt, text=ink},
  aibox/.style={draw=blue, line width=3.5pt, rounded corners=8pt, fill=skytint,
    align=flush center, inner sep=9pt, text=ink},
  ybox/.style={draw=yellowdeep, line width=3.5pt, rounded corners=8pt, fill=yellowtint,
    align=flush center, inner sep=9pt, text=ink},
  gbox/.style={draw=good, line width=3.5pt, rounded corners=8pt, fill=goodtint,
    align=flush center, inner sep=9pt, text=ink},
  obox/.style={draw=orange, line width=3.5pt, rounded corners=8pt, fill=orange!8,
    align=flush center, inner sep=9pt, text=ink},
  arr/.style={-{Stealth[length=18pt,width=17pt]}, line width=4pt, color=navy!70},
  dbl/.style={{Stealth[length=15pt,width=14pt]}-{Stealth[length=15pt,width=14pt]},
    line width=3.5pt, color=navy!70},
  lab/.style={font=\fontsize{20pt}{24pt}\selectfont, text=ink, align=left},
  elab/.style={font=\fontsize{19pt}{22pt}\selectfont, text=ink2},
  slab/.style={font=\fontsize{17pt}{20pt}\selectfont, text=ink2},   % 矢印が短いところのラベル
}

\pgfplotsset{
  gem/.style={
    axis lines=left, axis line style={ink2, line width=1.2pt},
    tick style={ink2, line width=1.2pt},
    tick label style={font=\fontsize{20pt}{22pt}\selectfont, color=ink},
    label style={font=\fontsize{20pt}{23pt}\selectfont, color=ink2},
    every axis plot/.append style={line width=3.5pt},
  },
}

\colorlet{orangetint}{orange!8}

\posttitle{AI の仕組みと評価}{2 / 2}

% ================================================================ 本文
\tikzposterlatexaffectionproofoff
\begin{document}
\maketitle[width=\paperwidth, linewidth=0pt, roundedcorners=0, innersep=7mm,
  titletoblockverticalspace=5mm]

\begin{columns}
% ================================================================ 左列
\column{0.5}

\block{\bt{割当 AI の仕組み}}{%
  \head{実績から「誰に任せるか」を学ぶ}
  {\centering
  \begin{tikzpicture}[font=\fontsize{21pt}{26pt}\selectfont]
    \tikzset{c/.style={text width=5.5cm, minimum height=3.8cm}}
    \newcommand{\st}[2]{{\color{blue}\bfseries\fontsize{32pt}{32pt}\selectfont #1}\ \ {\bfseries\fontsize{26pt}{30pt}\selectfont #2}\\[6pt]}
    \node[box, c] (a) {\st{1}{記録}所要時間\\本人の体感};
    \node[aibox, c, right=10mm of a] (b) {\st{2}{学習}人ごとの速さ\\＋\hl{自信の幅}};
    \node[aibox, c, right=10mm of b] (c) {\st{3}{選択}幅の中から\\くじを引く};
    \node[aibox, c, right=10mm of c] (d) {\st{4}{調整}手持ちで割る\\無資格者を除外};
    \foreach \i/\j in {a/b,b/c,c/d} \draw[arr] (\i) -- (\j);
    \draw[arr, rounded corners=12pt] (d.south) -- ++(0,-1.0) -| node[pos=0.25, below=2pt, elab]{作業が終わるたびにくり返す} (a.south);
  \end{tikzpicture}\par}
  \vspace{16pt}

  % ---- 手がかりと「うまくいったか」
  \begin{minipage}[t]{0.485\linewidth}\vspace{0pt}
    \tint{card}{%
      {\bfseries\fontsize{27pt}{32pt}\selectfont\color{navy}手がかり（タスクの特徴）\par}
      \vspace{8pt}
      {\bl
      \begin{tabular}{@{}>{\bfseries}l@{\hspace{10pt}}l@{}}
        難易度   & 1〜5 \\
        資格     & 溶接・クレーンなど \\
        機材     & 旋盤・プレスなど \\
        経験年数 & 勤続年数 \\
      \end{tabular}\par}
      \vspace{6pt}
      {\sm\color{ink2}資格・機材も見るので、\\「溶接が得意な人」も分かる\par}}
  \end{minipage}\hfill
  \begin{minipage}[t]{0.485\linewidth}\vspace{0pt}
    \tint{card}{%
      {\bfseries\fontsize{27pt}{32pt}\selectfont\color{navy}うまくいったか（報酬）\par}
      \vspace{10pt}
      {\centering\fontsize{25pt}{30pt}\selectfont
      $\text{体感}\times\dfrac{\text{いつもの 1 個の時間}}{\text{今回の 1 個の時間}}$\par}
      \vspace{10pt}
      {\sm
      \begin{tabular}{@{}l@{\hspace{14pt}}l@{\hspace{14pt}}l@{}}
        簡単 \textbf{1.5} & 普通 \textbf{1.0} & 難しい \textbf{0.5} \\
      \end{tabular}\par}
      \vspace{6pt}
      {\sm\color{ink2}数量で割り、1 個あたりで比べる\par}}
  \end{minipage}
  \vspace{18pt}

  % ---- 自信の幅とくじ
  \head{自信の幅から「くじ」を引く}
  \begin{minipage}[c]{0.50\linewidth}
    \begin{tikzpicture}
      \begin{axis}[gem, width=\linewidth, height=0.56\linewidth,
          xlabel={予想される速さ}, ylabel={確からしさ},
          xmin=0, xmax=2.2, ymin=0, ymax=5.6, ytick=\empty, xtick=\empty,
          samples=140, domain=0:2.2, clip=false]
        \addplot[color=navy, fill=navy, fill opacity=0.12]
          {exp(-(x-1.2)^2/(2*0.08^2))/(0.08*sqrt(2*pi))} \closedcycle;
        \addplot[color=orange, fill=orange, fill opacity=0.12]
          {exp(-(x-0.95)^2/(2*0.35^2))/(0.35*sqrt(2*pi))} \closedcycle;
        \draw[navy, line width=3.5pt, dashed] (axis cs:1.17,0) -- (axis cs:1.17,4.9);
        \draw[orange, line width=3.5pt, dashed] (axis cs:1.55,0) -- (axis cs:1.55,1.3);
        \node[lab, anchor=east, text=navy, align=right] at (axis cs:1.08,4.4)
          {ベテラン\\幅が狭い};
        \node[lab, anchor=west, text=orange, align=left] at (axis cs:1.42,2.4)
          {新人\\幅が広い};
      \end{axis}
    \end{tikzpicture}
    \capt{破線：くじで引いた値。大きい人に任せる}
  \end{minipage}\hfill
  \begin{minipage}[c]{0.47\linewidth}
    \begin{itemize}\fontsize{22pt}{29pt}\selectfont
      \item \hl{ベイズ線形回帰}：\\人ごとの速さと「自信の幅」
      \item \hl{Thompson Sampling}：\\幅の中からくじを引いて選ぶ
      \item 新人は幅が広く、\\ときどき大きな値を引く\\→ \hl{隠れた得意を発見}
    \end{itemize}
  \end{minipage}
  \vspace{18pt}

  % ---- 負荷分散と資格
  \begin{minipage}[t]{0.485\linewidth}\vspace{0pt}
    \tint{skytint}{%
      {\bfseries\fontsize{27pt}{32pt}\selectfont\color{blue}手持ちで割って偏りを防ぐ\par}
      \vspace{10pt}
      {\centering\fontsize{24pt}{30pt}\selectfont
      $\text{点数}=\dfrac{\text{くじの値}}{1+\dfrac{\text{手持ちの残り時間}}{\text{この仕事の見積もり}}}$\par}
      \vspace{8pt}
      {\sm 速いベテランでも、\\手持ちが多いと点数が下がる\par}}
  \end{minipage}\hfill
  \begin{minipage}[t]{0.485\linewidth}\vspace{0pt}
    \tint{orangetint}{%
      {\bfseries\fontsize{27pt}{32pt}\selectfont\color{orange}資格は学習させない\par}
      \vspace{10pt}
      {\fontsize{22pt}{29pt}\selectfont フォークリフト・溶接などは、\\AI を呼ぶ前に\hl{候補から除外}。\par}
      \vspace{8pt}
      {\sm 危険作業を「失敗から学ぶ」\\わけにはいかないため\par}}
  \end{minipage}
}

\block{\bt{2 通りの使い方と判断の根拠}}{%
  {\centering
  \begin{tikzpicture}[font=\fontsize{22pt}{27pt}\selectfont]
    \tikzset{c/.style={minimum height=1.9cm, inner sep=7pt}}
    \node[gbox, c, text width=5.2cm] (w) {{\bfseries\color{good} 現場でかざす}};
    \node[box, c, text width=15.4cm, right=16mm of w] (t) {その人に向く順に\hl{タスクを並べて}液晶へ};
    \node[box, c, text width=5.2cm, below=6mm of w] (t2) {{\bfseries 管理画面で\\「AI 割当」}};
    \node[gbox, c, text width=15.4cm, right=16mm of t2] (w2) {全員の点数を比べて\hl{担当者を選ぶ}};
    \draw[arr] (w) -- (t);
    \draw[arr] (t2) -- (w2);
  \end{tikzpicture}\par}
  \vspace{16pt}
  \head{「AI 割当」の根拠（実際の管理画面）}
  \begin{minipage}[c]{0.70\linewidth}
    \shot{0.97\linewidth}{admin_ai_why.png}\\[4pt]
    \capt{「製品F プレス抜き」の割当。\\実績 611 件（デモデータ）から学習}
  \end{minipage}\hfill
  \begin{minipage}[c]{0.28\linewidth}
    \raggedright\begin{itemize}\fontsize{23pt}{30pt}\selectfont
      \item 山口さんは\\\hl{手持ち 1.4 時間}\\→ 点数 0.67 に
      \item 手の空いた\\\hl{小林さん}を選択
      \item 理由が見えるので\\\hl{納得できる}
    \end{itemize}
  \end{minipage}
}

% ================================================================ 右列
\column{0.5}

\block{\bt{効果の検証（シミュレーション）}}{%
  {\bl 作業者 6 名（勤続 1〜30 年）で試した。\par
  各人に\hl{隠れた得意作業}があり、その作業は 4 割速い。\par
  AI には教えず、作業ログだけで見つけられるかを調べた。\par}
  \vspace{6pt}
  \capt{1 回に 6 件を割り当て × 150 回、20 試行の平均。割当 AI は実装をそのまま呼ぶ}
  \vspace{14pt}
  \begin{minipage}[c]{0.62\linewidth}
    \head{全員が終わるまでの時間}
    \begin{tikzpicture}
      \begin{axis}[gem, xbar, width=0.80\linewidth, height=0.50\linewidth,
          bar width=22pt, xmin=0, xmax=78, xtick=\empty, axis x line=none,
          ymin=0.4, ymax=4.6, ytick={1,...,4}, y dir=reverse,
          major y tick style={draw=none},
          yticklabels={経験年数順,ランダム,順番に回す,\textbf{Gemmba}},
          nodes near coords={\pgfmathprintnumber[fixed,precision=1,assume math mode=true]{\pgfplotspointmeta}\,分},
          nodes near coords align={horizontal},
          every node near coord/.append style={font=\fontsize{21pt}{21pt}\selectfont\bfseries, color=ink}]
        \addplot[fill=ink2!45, draw=none, bar shift=0pt] coordinates {(63.7,1) (44.4,2) (34.7,3)};
        \addplot[fill=blue, draw=none, bar shift=0pt] coordinates {(25.4,4)};
      \end{axis}
    \end{tikzpicture}
  \end{minipage}\hfill
  \begin{minipage}[c]{0.35\linewidth}
    {\centering
    {\fontsize{58pt}{58pt}\selectfont\bfseries\color{blue}−60\%}\\[6pt]
    {\fontsize{20pt}{24pt}\selectfont\color{ink2}経験年数順と比べて}\\[16pt]
    {\fontsize{58pt}{58pt}\selectfont\bfseries\color{blue}−27\%}\\[6pt]
    {\fontsize{20pt}{24pt}\selectfont\color{ink2}順番に回す割当と比べて}\par}
  \end{minipage}
  \vspace{16pt}

  \head{使うほど早くなる（学習の進み方）}
  {\centering
  \begin{tikzpicture}
    \begin{axis}[gem, width=0.98\linewidth, height=0.40\linewidth,
        xlabel={割当の回数}, ylabel={全員が終わるまで [分]},
        xmin=0, xmax=150, ymin=15, ymax=72, xtick={0,50,100,150}, ytick={20,40,60},
        unbounded coords=jump, clip=false,
        legend style={draw=none, fill=none, font=\fontsize{19pt}{22pt}\selectfont,
          at={(1.0,1.0)}, anchor=north east, legend columns=3,
          /tikz/every even column/.append style={column sep=14pt}},
        legend cell align=left]
      \pgfplotstableread[col sep=comma]{../../resume/poster_curve.csv}\curve
      \addplot[ink2!45, line width=3pt] table[x=round, y=veteran_span]{\curve};
      \addlegendentry{経験年数順}
      \addplot[ink2!80, line width=3pt, dashed] table[x=round, y=random_span]{\curve};
      \addlegendentry{ランダム}
      \addplot[navy!55, line width=3pt, dotted] table[x=round, y=roundrobin_span]{\curve};
      \addlegendentry{順番に回す}
      \addplot[blue, line width=4.5pt] table[x=round, y=ts_load_span]{\curve};
      \addlegendentry{\textbf{Gemmba}}
    \end{axis}
  \end{tikzpicture}\par}
  \capt{はじめは手探りでも、ログがたまるほど短くなる（20 試行の平均）}
  \vspace{14pt}

  \tint{yellowtint}{%
    {\bl\bfseries 分かったこと：「一番向く人」と「全体が早く終わる割当」は違う\par}
    \vspace{6pt}
    {\fontsize{23pt}{28pt}\selectfont
    \renewcommand{\arraystretch}{1.25}
    \begin{tabular*}{\linewidth}{@{\extracolsep{\stretch{1}}}l r r r@{}}
      \hline
      & 向いている度\textsuperscript{*} & 1 人への集中 & 完了時間 \\
      \hline
      手持ちで割らない & \textbf{1.03} & 42\% & 32.4 分 \\
      \textbf{\color{blue}手持ちで割る} & 1.16 & \textbf{27\%} & \textbf{\color{blue}25.4 分} \\
      \hline
    \end{tabular*}\par}
    \vspace{4pt}
    {\fontsize{18pt}{22pt}\selectfont\color{ink2}* 任せた人の時間 ÷ 最も向く人の時間（1 が最良）。学習後（51 回目以降）で計測\par}
    \vspace{6pt}
    {\fontsize{23pt}{29pt}\selectfont ログだけで隠れた得意を見つけ、ほぼ最適な人を選べる。\par
    ただし仕事が偏るので、\hl{手持ちで割る工夫}が要る。\par}}
}

\block{\bt{音声 AI：話すだけでタスクを登録}}{%
  {\centering
  \begin{tikzpicture}[font=\fontsize{21pt}{26pt}\selectfont]
    \tikzset{c/.style={text width=5.5cm, minimum height=2.7cm}}
    \node[box, c] (v1) {{\bfseries 話す}\\[4pt]{\sm\color{ink2}ボタンを押して}};
    \node[ybox, c, right=10mm of v1] (v2) {{\bfseries 文字にする}\\[4pt]Whisper};
    \node[ybox, c, right=10mm of v2] (v3) {{\bfseries 意味を読む}\\[4pt]Gemma 3（LLM）};
    \node[box, c, right=10mm of v3] (v4) {{\bfseries 検査して登録}\\[4pt]{\sm\color{ink2}資格を確認}};
    \foreach \i/\j in {v1/v2,v2/v3,v3/v4} \draw[arr] (\i) -- (\j);
  \end{tikzpicture}\par}
  \vspace{14pt}
  \begin{minipage}[t]{0.31\linewidth}\vspace{0pt}\centering
    \photo{uc4_rec.png}{\linewidth}{0.64\linewidth}{}\\[4pt]
    \capt{ボタンを押している\\間だけ録音}
  \end{minipage}\hfill
  \begin{minipage}[t]{0.31\linewidth}\vspace{0pt}\centering
    \photo{uc4_done.png}{\linewidth}{0.64\linewidth}{}\\[4pt]
    \capt{登録できたら\\液晶に表示}
  \end{minipage}\hfill
  \begin{minipage}[t]{0.345\linewidth}\vspace{0pt}
    \begin{itemize}\fontsize{22pt}{28pt}\selectfont
      \item \hl{現場の PC で処理}\\情報を外に出さない
      \item 気づいた人が\\その場で登録\\→ \hl{忘れを防ぐ}
    \end{itemize}
  \end{minipage}
  \vspace{12pt}
  {\fontsize{20pt}{25pt}\selectfont
  \renewcommand{\arraystretch}{1.3}
  \begin{tabular*}{\linewidth}{@{\extracolsep{\stretch{1}}}l l c c c@{}}
    \hline
    話した言葉 & タスク名 & 重要度 & 数量 & 期限 \\
    \hline
    「旋盤の切粉清掃を至急お願いします」 & 旋盤の切粉清掃 & \textcolor{red!70!black}{\textbf{至急}} & 1 & -- \\
    「製品 C の外形加工を 20 個、明日まで…」 & 製品Cの外形加工 & \textbf{高} & 20 & 翌日 \\
    \hline
  \end{tabular*}\par}
  \capt{仮想モジュールから送った言葉と、Gemma 3 が実際に作ったタスク}
}

\end{columns}

% ---------------------------------------------------------------- ユースケース（全幅）
\block{\bt{ユースケース}}{%
  % 1 行：場面（太字）と、そのときの操作の流れ
  \newcommand{\ucrow}[2]{%
    \begin{minipage}[t]{0.25\linewidth}{\bfseries\fontsize{24pt}{29pt}\selectfont #1}\end{minipage}\hfill
    \begin{minipage}[t]{0.74\linewidth}{\fontsize{21pt}{29pt}\selectfont #2}\end{minipage}\par
    \vspace{6pt}}
  \newcommand{\ucbox}[4]{% 地の色, 文字色, 見出し, 行
    \begin{minipage}[t]{0.488\linewidth}\vspace{0pt}
      \tint{#1}{%
        {\fontsize{32pt}{38pt}\selectfont\bfseries\color{#2}#3}\par
        \vspace{8pt}
        \def\ucto{\ \textcolor{#2}{→}\ }#4}
    \end{minipage}}
  \ucbox{goodtint}{good}{従業員：}{%
    \ucrow{仕事をする}{腕輪をかざす\ucto 提案を承認\ucto 終わったら再度かざす}
    \ucrow{急ぎの仕事}{かざすと至急の仕事へ誘導\ucto 移動先にも通知}
    \ucrow{仕事に気づく}{ボタンを押して話す\ucto AI がタスクに登録}}\hfill
  \ucbox{skytint}{blue}{管理者：}{%
    \ucrow{仕事を割り振る}{タスクを登録\ucto 「AI 割当」\ucto 担当者が決まる}
    \ucrow{現場を確かめる}{機材の使用中・空きと、誰が何をしているかを確認}
    \ucrow{人を育てる}{ログで得意・不得意を把握。新人にも試しに任せる}}
}

\end{document}
